\section{How TimesFM 3 works}
\label{sec:mechanism}

TimesFM~3 is described publicly by a blog post and a model card \citep{jain2026timesfm3,timesfm2026hf}; there is no technical report. We
therefore read the inference code of the public package (\texttt{timesfm} 3.0.2, PyTorch backend, commit \texttt{e51928e}) line by
line, together with the checkpoint configuration (\texttt{google/timesfm-3.0-pytorch}, revision \texttt{43046b85}). The
description below is limited to what the code shows; statements about training are inferences and are marked as such.
Figure~\ref{fig:architecture} summarises the inference path.

\begin{figure}[!htbp]
\centering
\begin{tikzpicture}[
  font=\footnotesize, node distance=0.22cm,
  blk/.style={draw=vzaxis, rounded corners=3pt, align=center, inner sep=3pt, text width=3.35cm, minimum height=1.05cm,
    execute at begin node={\hyphenpenalty=10000\exhyphenpenalty=10000}},
  fixed/.style={blk, fill=vzorangewash},
  learned/.style={blk, fill=vzbluewash},
  arr/.style={-{Stealth[length=1.8mm]}, thick, vzink2}
]
\node[fixed] (in) {\textbf{API preprocessing}\\ interpolate gaps, keep last\\ 15{,}360 points, round $H$ to 64};
\node[fixed, right=of in] (det) {\textbf{Linear detrending gate}\\ OLS line per variate, used\\ only if $R^2 > 0.75$};
\node[learned, right=of det] (tok) {\textbf{Patch tokens}\\ 32 points, causal running\\ RevIN; 192 input features};
\node[learned, right=of tok] (tr) {\textbf{20 mixing layers}\\ causal attention over time,\\ then over variates; FFN};
\node[learned, below=0.75cm of tr] (head) {\textbf{Quantile head}\\ each token: next 64 points\\ $\times$ 9 quantiles};
\node[fixed, left=of head] (cpm) {\textbf{CPM iterative RevIN}\\ statistics rolled forward\\ with its own median};
\node[fixed, left=of cpm] (st) {\textbf{Stitching}\\ 64-point windows, hop 32,\\ linear cross-fade};
\node[fixed, left=of st] (out) {\textbf{Output}\\ add back the line; sort\\ quantiles; median = point};
\draw[arr] (in) -- (det); \draw[arr] (det) -- (tok); \draw[arr] (tok) -- (tr); \draw[arr] (tr) -- (head);
\draw[arr] (head) -- (cpm); \draw[arr] (cpm) -- (st); \draw[arr] (st) -- (out);
\node[below=0.15cm of out.south west, anchor=north west, font=\scriptsize, text=black!70, text width=15.2cm]
  {\colorbox{vzbluewash}{\strut learned} \colorbox{vzorangewash}{\strut fixed computation (no parameters)}\quad
   One forward pass: horizon positions enter as masked tokens (CPM); there is no autoregressive loop.
   For $H \le 64$ the whole forecast comes from the last context token.};
\end{tikzpicture}
\caption{Inference path of TimesFM~3 as implemented in the public PyTorch code (\texttt{timesfm} 3.0.2), read line by line
(Section~\ref{sec:mechanism}). Shaded boxes are deterministic steps without parameters; four of them (detrending, CPM refinement,
stitching, quantile sorting) can shape a forecast, some only conditionally (the detrending gate; CPM and stitching only beyond 64
steps), but are not part of the network.}
\label{fig:architecture}
\end{figure}


\paragraph{Tokens and normalisation.} A series is cut into patches of 32 points aligned to its last observation; shorter batches
are left-padded with masked zeros, and the public API keeps at most the last 15{,}360 points. Each patch is normalised with the
mean and standard deviation of all observed points up to and including that patch (a causal, ``running'' RevIN), and the token
receives the 32 normalised values, a 64-point look-ahead channel and two mask channels (192 features), mapped to the model width
1{,}280 by a residual block. The look-ahead channel is non-empty only for covariates whose future values are known; for targets and
past-only covariates it is masked. Interior gaps are filled by linear interpolation before the model sees them: the network can take
point masks, but the forecasting API does not pass them, so missingness never reaches the model.

\paragraph{Transformer.} Twenty layers alternate, inside each layer, causal attention along the patch axis (rotary positions) and
full attention across variates at the same patch index, followed by a feed-forward block with no expansion (1{,}280 to 1{,}280);
the parameter count implied by the code shapes is about 330.7 million, as advertised. The checkpoint configuration sets
\texttt{use\_rope\_var = false}: attention across variates has no positional signal, so the model is equivariant to permutations of
the variates. A ``past-only covariate'' is encoded exactly as an extra target and is forecast and then discarded; inside the network
the two roles are indistinguishable. The limit of 32 variates mentioned in the documentation is not enforced by the model or by the
forecasting API; only the evaluator class used by Google's benchmark runners splits larger problems into chunks of 32 and randomly
subsamples covariates beyond that number.

\paragraph{Horizon and decoding.} Forecasting is a single forward pass. Every token predicts the next 64 points for nine quantiles
(0.1 to 0.9); the horizon is represented by masked tokens appended to the context (continuous patch masking, CPM). Overlapping
64-point windows of consecutive tokens, with a hop of 32, are cross-faded linearly (``stitching''). Because the statistics used to
de-normalise a horizon token cannot be observed, an iterative routine rolls them forward by treating the model's own median forecast
as if it were observed. Two consequences follow from the code. For horizons up to 64 steps, the entire forecast is produced by the
last context token, and neither the masked tokens nor the refinement play any role. With causal attention and no known future
covariates, the forecast of the first $h$ steps does not depend on the requested horizon.

\paragraph{Fixed computations around the network.} Four steps shape every forecast but have no parameters. (i) \emph{Linear
detrending:} each variate gets an ordinary least-squares line over its context, and the residuals are used instead of the series
when $\mathrm{std}(\text{residual}) < 0.5\,\mathrm{std}(\text{series})$, i.e.\ when the line explains more than 75\% of the
variance; the extrapolated line is then added back to all nine quantiles. The gate is discontinuous: a small change of the context
can switch it on or off. (ii) The CPM refinement described above. (iii) Stitching. (iv) Quantile sorting: the nine heads are
independent linear outputs, so they are sorted, and the point forecast is the fifth sorted value. Two further options are off by
default in the forecasting API but on in all of Google's benchmark runners: \emph{symmetric averaging}, which forecasts $x$ and $-x$
and averages the reflected quantiles (doubling the compute), and \emph{make positive}, which clips the quantiles of non-negative
series at zero. We call the package defaults \texttt{tfm3} and the runners' flags \texttt{tfm3-eval}.

\paragraph{Properties worth testing.} The design implies per-variate equivariance to positive affine transformations (each variate
is normalised by its own location and scale), equivariance to permutations of variates, no learned mechanism for missing data, and a
hard-coded linear extrapolation that may be harmful on saturating or regime-switching trends. Section~\ref{sec:controlled} tests the
last point (H3) and the ablations in Section~\ref{sec:ablations} switch each fixed step off in turn. The full reading of the code,
with line references, is released as \texttt{reports/mechanism.md}.
